\documentclass[10pt,a4paper]{amsart}
\usepackage[margin=19mm]{geometry}
\usepackage{amsmath,amssymb,mathtools}
\usepackage[scr=rsfs,cal=euler]{mathalfa}
\usepackage{xfrac}
\usepackage[unicode,hidelinks,bookmarks=false]{hyperref}
\newcommand{\ie}{i.e.\ }
\newcommand{\cf}{cf.\ }
\newcommand{\resp}{respectively\ }
\newcommand{\Subsection}{\subsection}
\providecommand{\nobreakdash}{\nobreak\hbox{-}}
\setlength{\emergencystretch}{2em}
\multlinegap0pt
\usepackage{amssymb}
\usepackage{xcolor}
\usepackage{enumerate}
\usepackage{float}
\usepackage{braket}
\usepackage{algorithm}\usepackage{algpseudocode}
\usepackage{mathtools}
\usepackage{bm}
\newtheorem{theorem}{Theorem}
\title{\textbf{Flagged Extensions and Numerical Simulations for Quantum Channel Capacity: Bridging Theory and Computation}}
\author{Vahid Nourozi}
\date{Source: arXiv:2506.03429v2 (CC BY 4.0)}
\begin{document}
\maketitle

\noindent\emph{Source and licence.} \textbf{Flagged Extensions and Numerical Simulations for Quantum Channel Capacity: Bridging Theory and Computation}. Version arXiv:2506.03429v2 (CC BY 4.0), distributed by arXiv under the Creative Commons Attribution 4.0 International licence, http://creativecommons.org/licenses/by/4.0/, as recorded in the arXiv licence metadata for this exact version and shown on the abstract page https://arxiv.org/abs/2506.03429v2. Full licence notice: \texttt{licenses/arxiv-2506.03429-CC-BY-4.0.txt}.\par\smallskip

\noindent\textit{Editorial scope.} Counted unit: the article's theorem th22 (source0.tex lines 149--160). The article's complete original proof of it is delivered with it (source0.tex lines 163--193, one \texttt{proof} environment of 31 lines). No mathematics is written, repaired or re-derived: the statement and the proof are the article's own lines, in the article's own notation, and no retained line is substituted or re-flowed.  No further source-local context is needed: every cross-reference of every retained line resolves inside the two delivered spans.  Nothing was omitted from any retained span, so the package carries no omission record.  No retained line contains a citation, so the wrapper carries no bibliography and no bibliography entry.  No retained line contains a figure environment, an image-inclusion command or drawing code, so no image is delivered and the package declares no asset dependency.  Editorial formatting only, in addition: an AMS \texttt{amsart} preamble that restates the article's own declarations and macros used by the retained lines, namely the environments \texttt{theorem} on separate counters, together with the standard packages those lines need.  The retained environments print in this document as Theorem 1 in order of appearance, whereas the article prints the same items with the numbering of its own full document.  The complete native source of the article as distributed in the arXiv e-print for this exact version is bundled unchanged as \texttt{sources/179265/source0.tex}.
\begin{theorem}[Flagged extension is degradable when the noise branch is degradable]\label{th22}
Let $N = p\,U + (1-p)\,M$ be a convex mixture of a unitary channel $U(\rho)=U\rho U^\dagger$ and an arbitrary CPTP map $M$, with $0<p\le 1$. Define the flagged extension
\[
\widetilde N(\rho)= p\,|0\rangle\!\langle 0|_F\otimes U(\rho) + (1-p)\,|1\rangle\!\langle 1|_F\otimes M(\rho).
\]
If $M$ is degradable, then $\widetilde N$ is degradable. Consequently,
\[
Q(N)\le Q(\widetilde N)= I_c(\rho_\star,\widetilde N),\qquad 
P(N)\le P(\widetilde N),
\]
and both $Q(\widetilde N)$ and $P(\widetilde N)$ are single–letter.
\end{theorem}

\begin{proof}
Let $V_U:\mathcal H_A\!\to\!\mathcal H_B\!\otimes\!\mathcal H_{E_U}$ and $V_M:\mathcal H_A\!\to\!\mathcal H_B\!\otimes\!\mathcal H_{E_M}$ be Stinespring isometries for $\mathcal U$ and $\mathcal M$, so that
$\mathcal U(\rho)=\mathrm{Tr}_{E_U}\!\left[V_U\rho V_U^\dagger\right]$ and
$\mathcal U^c(\rho)=\mathrm{Tr}_{B}\!\left[V_U\rho V_U^\dagger\right]$ (similarly for $\mathcal M,\mathcal M^c$).
A Stinespring isometry for $\widetilde{\mathcal N}$ is
\[
W:\mathcal H_A \to \mathcal H_F\!\otimes\!\mathcal H_B\!\otimes\!\mathcal H_E,\qquad
W=\sqrt{p}\,|0\rangle_F\!\otimes V_U \;\oplus\; \sqrt{1-p}\,|1\rangle_F\!\otimes V_M,
\]
where I identify $\mathcal H_E\simeq \mathcal H_{E_U}\oplus \mathcal H_{E_M}$ and $\mathcal H_F=\mathrm{span}\{|0\rangle,|1\rangle\}$ is the flag.
Tracing out $E$ yields $\widetilde{\mathcal N}$; tracing out $B$ yields $\widetilde{\mathcal N}^c$ as in the statement.

Define $\mathsf D:\mathcal B(\mathcal H_F\!\otimes\!\mathcal H_B)\!\to\!\mathcal B(\mathcal H_F\!\otimes\!\mathcal H_E)$ by
\[
\mathsf D(\sigma)= |0\rangle\!\langle 0|_F\!\otimes\!\mathsf D_0(\langle 0|\sigma|0\rangle_F)\;+\;|1\rangle\!\langle 1|_F\!\otimes\!\mathsf D_1(\langle 1|\sigma|1\rangle_F),
\]
with branch maps
\[
\mathsf D_0(\tau_B) \;=\; \mathcal U^c(\rho_\star)\quad(\text{a fixed state on }E_U),\qquad
\mathsf D_1(\tau_B) \;=\; \Delta(\tau_B).
\]
Here $\rho_\star$ is any fixed state on $A$; since $\mathcal U$ is unitary, $\mathcal U^c(\cdot)$ is constant (isometric output on $E_U$ independent of the input to $B$), so $\mathsf D_0$ is CPTP and does not depend on $\tau_B$. By assumption, $\Delta$ is a CPTP degrading map for $\mathcal M$, hence $\mathsf D_1$ is CPTP. Therefore $\mathsf D$ is CPTP.

Applying $\mathsf D$ to $\widetilde{\mathcal N}(\rho)=p\,|0\rangle\!\langle 0|\!\otimes\!\mathcal U(\rho)+(1-p)\,|1\rangle\!\langle 1|\!\otimes\!\mathcal M(\rho)$ gives
\[
\mathsf D\!\big(\widetilde{\mathcal N}(\rho)\big)=
p\,|0\rangle\!\langle 0|\!\otimes\!\mathcal U^c(\rho_\star)\;+\;(1-p)\,|1\rangle\!\langle 1|\!\otimes\!\Delta\big(\mathcal M(\rho)\big).
\]
Because $\mathcal U^c$ is independent of the $B$-output and $\Delta\circ\mathcal M=\mathcal M^c$, the right-hand side equals
$p\,|0\rangle\!\langle 0|\!\otimes\!\mathcal U^c(\rho)\;+\;(1-p)\,|1\rangle\!\langle 1|\!\otimes\!\mathcal M^c(\rho)=\widetilde{\mathcal N}^c(\rho)$ up to the harmless choice of $\rho_\star$ for the unitary branch (any fixed choice matches the complementary output because $\mathcal U^c$ is constant). Hence $\widetilde{\mathcal N}$ is degradable.
\end{proof}

\end{document}
